\chapter{Thực và Ảo}
\markboth{Thực và Ảo}{Real and Illusion}

\section{Một người sống rất đẹp trên mạng nhưng ngoài đời kiệt sức.}
\begin{truyen}{Một người sống rất đẹp trên mạng nhưng ngoài đời kiệt sức.}{A polished image can hide a collapsing life.}
\chuhoa{A}{nh biết cách chụp góc đẹp, nói câu đẹp và tạo cảm giác cuộc đời mình rất ổn. Nhưng ngoài màn hình, anh ngủ ít, lo nhiều và không thật sự vui. Ảo ảnh càng chỉn chu, đời thật càng dễ bị bỏ quên.}
\textit{(He knew how to take beautiful photos, write attractive lines, and make life look stable. Outside the screen, he slept little, worried much, and was not truly happy. The more polished the illusion, the easier real life is neglected.)}
\end{truyen}
\begin{baihoc}
Thứ trông đẹp chưa chắc đang sống khỏe.
\textit{(What looks beautiful is not always living well.)}
\end{baihoc}
\begin{ghinhoanh}
\textbf{Short English to remember:}

A polished image can hide a tired life.
\end{ghinhoanh}
\begin{tuhocnhanh}
1. Em có đang chăm hình ảnh nhiều hơn chăm đời sống thật không? \textit{(Are you caring more for image than for real life?)}

2. Điều gì trong đời thật của em cần được cứu trước tiên? \textit{(What in your real life needs help first?)}
\end{tuhocnhanh}
\ngancach

\section{Một học sinh tưởng mình hiểu bài vì xem lời giải rất nhiều.}
\begin{truyen}{Một học sinh tưởng mình hiểu bài vì xem lời giải rất nhiều.}{Recognition is not mastery.}
\chuhoa{E}{m đọc qua thì thấy quen, nhìn lại công thức thì thấy có vẻ hiểu. Nhưng đến lúc tự làm, tay lại đứng lại. Em mới biết cảm giác quen mặt với kiến thức không đồng nghĩa với việc thật sự hiểu và làm được.}
\textit{(He felt familiar with the lesson after seeing many worked examples. But when he tried alone, his hands stopped. He learned that recognizing knowledge is not the same as truly understanding and applying it.)}
\end{truyen}
\begin{baihoc}
Cảm giác biết là một trong những ảo tưởng dễ chịu nhất của việc học.
\textit{(The feeling of knowing is one of learning's most comfortable illusions.)}
\end{baihoc}
\begin{ghinhoanh}
\textbf{Short English to remember:}

Recognition is not mastery.
\end{ghinhoanh}
\begin{tuhocnhanh}
1. Em có đang nhầm quen với hiểu không? \textit{(Are you confusing familiarity with understanding?)}

2. Phần nào em phải tự làm mới biết thật sự đã hiểu? \textit{(What must you do on your own to know you truly understand?)}
\end{tuhocnhanh}
\ngancach

\section{Một người tưởng mình được yêu vì luôn được chiều.}
\begin{truyen}{Một người tưởng mình được yêu vì luôn được chiều.}{Indulgence is not always love.}
\chuhoa{C}{ó lúc người ta chiều mình vì sợ mất mình, vì ngại xung đột, hoặc vì mỏi mệt chứ không hẳn vì yêu sâu. Khi mối quan hệ đứt ra, cô mới hiểu được yêu không chỉ là được nhường, mà là được cùng nhau lớn lên.}
\textit{(Sometimes people indulge us because they fear losing us, avoid conflict, or feel tired, not because they love deeply. When the relationship ended, she learned that being loved is not only being indulged, but growing together.)}
\end{truyen}
\begin{baihoc}
Thứ làm ta dễ chịu chưa chắc là thứ làm ta được yêu đúng nghĩa.
\textit{(What makes us comfortable is not always what means we are truly loved.)}
\end{baihoc}
\begin{ghinhoanh}
\textbf{Short English to remember:}

Indulgence is not always love.
\end{ghinhoanh}
\begin{tuhocnhanh}
1. Em có đang nhầm sự chiều chuộng với tình yêu không? \textit{(Are you confusing indulgence with love?)}

2. Tình yêu đúng nghĩa cần thêm điều gì ngoài dễ chịu? \textit{(What does real love need beyond comfort?)}
\end{tuhocnhanh}
\ngancach

\section{Một người tưởng mình tự do vì không bị ai nhắc nhở.}
\begin{truyen}{Một người tưởng mình tự do vì không bị ai nhắc nhở.}{Lack of rules is not the same as freedom.}
\chuhoa{C}{ậu muốn không ai quản, không ai yêu cầu, không ai kiểm tra. Khi sống như vậy một thời gian, cậu không thấy tự do hơn mà thấy đời rời rạc và thiếu hướng đi. Không bị ràng buộc chưa chắc là tự do; có khi chỉ là đang trôi.}
\textit{(He wanted no one to supervise, demand, or check on him. After living like that for a while, he felt not freer but more scattered and directionless. The absence of structure is not always freedom; sometimes it is drifting.)}
\end{truyen}
\begin{baihoc}
Tự do thật không phải muốn gì làm nấy, mà là biết tự giữ mình khỏi tan ra.
\textit{(Real freedom is not doing whatever you want, but being able to keep yourself from falling apart.)}
\end{baihoc}
\begin{ghinhoanh}
\textbf{Short English to remember:}

Lack of rules is not the same as freedom.
\end{ghinhoanh}
\begin{tuhocnhanh}
1. Em có đang gọi sự buông thả là tự do không? \textit{(Are you calling looseness freedom?)}

2. Kỷ luật nào giúp em tự do hơn thật sự? \textit{(What discipline actually makes you freer?)}
\end{tuhocnhanh}
\ngancach

\section{Một người tưởng mình sống rất thật nhưng chỉ đang nói mọi thứ không lọc.}
\begin{truyen}{Một người tưởng mình sống rất thật nhưng chỉ đang nói mọi thứ không lọc.}{Brutal honesty is not always true maturity.}
\chuhoa{A}{nh tự hào vì mình thẳng, nghĩ gì nói nấy. Nhưng nhiều lần lời thẳng ấy làm người khác đau mà không tạo ra điều gì tốt hơn. Về sau anh mới hiểu sống thật không có nghĩa là bỏ đi sự tinh tế và lòng trắc ẩn.}
\textit{(He was proud of being blunt and speaking whatever came to mind. But many times that bluntness only hurt others without creating anything better. Later he understood that authenticity does not mean abandoning tact and compassion.)}
\end{truyen}
\begin{baihoc}
Sống thật là thật với giá trị của mình, không phải thật với mọi cơn bốc đồng của miệng lưỡi.
\textit{(Living truthfully means being true to your values, not to every impulse of your tongue.)}
\end{baihoc}
\begin{ghinhoanh}
\textbf{Short English to remember:}

Bluntness is not the same as maturity.
\end{ghinhoanh}
\begin{tuhocnhanh}
1. Em có đang gọi sự thiếu tinh tế là sống thật không? \textit{(Are you calling a lack of tact authenticity?)}

2. Em muốn lời thật của mình đi cùng điều gì hơn? \textit{(What do you want your honest words to travel with?)}
\end{tuhocnhanh}
\ngancach